\documentclass{article}

% Pakete
\usepackage[utf8]{inputenc}
\usepackage{bookmark}
\usepackage[T1]{fontenc}
\usepackage[ngerman]{babel}
\usepackage{geometry}
\usepackage{setspace}
\usepackage{hyperref}

\geometry{margin=2.5cm}
\onehalfspacing

\title{%
    Exposé der Bachelorarbeit\\[1em]
    \large Range Minimum Queries in multiple dimensions
}

\author{%
    Name: Samra Abd El Wahab\\
    Matrikelnummer: 7214480\\
    Studiengang: Informatik B.Sc.\\
    Betreuer: Dr. Jonathan Rollin
}

\date{\today}

\begin{document}

\maketitle

\section{Einleitung und Problemstellung}
% Beschreibe hier:
% - Worum geht es?
% - Warum ist das Thema relevant?
% - In welchem Kontext steht es?
Was ist die Problemstellung, welche Varianten gibt es, wie funktioniert der Algo?

\section{Ziel der Arbeit}
% Was willst du erreichen?
% Was ist dein konkreter Beitrag?

\section{Forschungsfragen}
% Formuliere 2–4 konkrete Fragen
\begin{itemize}
    \item Frage 1
    \item Frage 2
\end{itemize}

\section{Stand der Forschung}
% Kurzer Überblick über bestehende Ansätze
% (keine Details, nur Einordnung)

\section{Methodik}
% Wie gehst du vor?
% Beispiele:
% - Analyse
% - Implementierung
% - Vergleich

\section{Geplanter Aufbau der Arbeit}
\begin{enumerate}
    \item Introduction
    \item Fundamentals 
        \begin{enumerate}
            \item Introduce 1D-RMQ using LCA as well as the direct approach
            \item Problem defintion and further definitions
            \item Show the limitation of using cartesian trees to solve 2D-RMQ problem
        \end{enumerate}
    \item Naive and Previous Approaches
        \begin{enumerate}
        \item A naive approach and alaysis of it
        \item Papers that come before this
        \end{enumerate}
    \item $\langle O(n^d), O(1) \rangle$ Approach
        \begin{enumerate}
            \item Yuan's paper: Detailed shwoing of methodolgy and proofs
            \item Thus we have a linear preprocessing and constant querying time algorithm, similar to the 1D-RMQ case
            \item Analysis
        \end{enumerate}
    \item Extension to Symmetric Queries
        \begin{enumerate}
            \item Problem Reformulation 
            \item Challenges
            \item Proposed solution
            \item Analysis and conclusion of solution, whether it brings any improvements
        \end{enumerate}
    \item Further Developments and Applications
    \item Discussion
    \item Conclusion
\end{enumerate}

\section{Plan}
\begin{tabular}{ll}
Week 1 & Research \\
Week 2--3 & Fundamentals \\
Week 4--5 & Yuan \\
Week 6--7 & Extension to Symmetric Queries \\
Week 8--10 & Further developments, discussion and conclusion \\
Week 10--12 & Korrektur, feinschliff \\
\end{tabular}

\section{Literatur}
% Erste Quellen
\begin{itemize}
    \item Autor, Titel, Jahr
    \item Autor, Titel, Jahr
\end{itemize}

\end{document}